% !TEX root = main.tex
% Main text: abstract through conclusion. Proofs are in the appendix; numerical checks in Paper2/checks/.
\begin{abstract}
Two training distributions can agree on every local statistic, every Bayes predictor and, up to a vanishing total-variation distance, on everything else,
yet make gradient training of the same attention model polynomially slower. We prove this for a one-block attention--ReLU Transformer whose
reduced dynamics are derived from its matrices. The mechanism is credit cancellation: attention reads every prefix position, so correlations of order
$\tfrac{1}{N}$ between the source and the other positions, invisible one at a time, add up and cancel the gradient credit reaching the rule.
With independent context the routing time is $\Theta(N\log(1+\varepsilon^{-2}))$ from initialization scale $\varepsilon$; with cancelled credit it is $\Theta(\tfrac{N}{\varepsilon})$,
and over-cancellation stops learning entirely below an exact initialization barrier. The cancellation can be made distributional, so that no function
of what uniform attention reads carries rule credit; a single earlier step that displays the negation of the current state suffices.
Corrupted chain-of-thought supervision is our motivating case: a causal error process with a constant fraction of wrong steps produces the delay at every supervised position,
while independent errors only rescale time by $\tfrac{1}{1-2\eta}$.
The bottleneck is fragile in a useful way: mixing a fraction $\lambda$ of traces without cancellation gives routing time
$\Theta\bigl(N\min\{\tfrac{1}{\varepsilon},\tfrac{1+\log(1+\lambda/\varepsilon)}{\lambda}\}\bigr)$, so $\lambda$ of the order of $\varepsilon$ already removes it.
Exact population experiments check every rate. In a block whose every matrix entry trains, started from the specified directions with or without noise,
the same three regimes appear; from fully random initialization they do not separate in our runs.
\end{abstract}

% ===================================================================
\section{Introduction}
\label{sec:intro}

When can two data distributions that look the same to every local test train a network at very different speeds?
We give an exact answer for attention. An attention head that must learn where to look reads all $N$ prefix positions at once.
If each of them correlates with the target-relevant position by only $-\tfrac{1}{N-1}$, no single statistic notices, but their sum cancels the
gradient credit that would teach the network its rule. The network then sits near its initialization for a time of order $\tfrac{1}{\varepsilon}$
instead of $\log\tfrac{1}{\varepsilon}$.

\paragraph{A saddle made by the data.}
Slow escape from small initialization is familiar from deep linear networks, where a product of small factors creates a saddle
\citep{saxe2014exact,shamir2019exponential}. The saddle here has a different origin. Our model has the same depth under both laws and the same
source--target correlation; what differs is how the other positions correlate with the source. The cancellation turns the linear growth of a
conserved primitive into quadratic growth, and that changes the escape time from logarithmic to inverse in $\varepsilon$.
Weak-gradient analyses \citep{shalevshwartz2017failures,abbe2025parities} relate slow learning to weak target correlation; here the target correlation is fixed
and only the sum over attended positions cancels.

\paragraph{Corrupted chain-of-thought as the motivating case.}
Reasoning traces used for supervision can contain incorrect intermediate steps \citep{havrilla2024noise,zhou2024noisyrationales},
and a model can still learn the correct rule when it remains the most likely next step \citep{natarajan2013noisy,patrini2017loss}.
For modulo-two counting the clean step is $x_{i+1}=b_{i+1}x_i$; a corrupted trace shows $x_{i+1}=b_{i+1}x_i\xi_{i+1}$ with $\Pr(\xi_{i+1}=-1)=\eta$.
At uniform attention the rule gradient factorizes into reliability and contextual credit:
\[
 -\pd{\mathcal L}{r_1}
 =\frac{r_2}{2N}
 \underbrace{(1-2\eta)}_{\text{label reliability }\gamma}
 \underbrace{\left(1+\sum_{j=1}^{N-1}\E[x_Nx_j]\right)}_{\text{contextual credit }\kappa}.
\]
Independent errors lower $\gamma$, which only rescales time. Errors that make earlier steps anti-correlate with the current state lower $\kappa$,
which changes the order of the routing time. A single earlier step that displays the negation of the current state is enough.

\begin{figure}[t]
\centering
\resizebox{\linewidth}{!}{% Conceptual figure: what uniform attention reads under the two matched laws.
\begin{tikzpicture}[>=Stealth, font=\small,
  src/.style={rectangle, rounded corners=2pt, draw=royalazure!80!black, fill=royalazure!14, minimum width=9mm, minimum height=6mm, inner sep=1pt},
  dec/.style={rectangle, rounded corners=2pt, draw=black!45, fill=black!5, minimum width=9mm, minimum height=6mm, inner sep=1pt},
  anti/.style={rectangle, rounded corners=2pt, draw=red!70!black, fill=red!10, minimum width=9mm, minimum height=6mm, inner sep=1pt},
  qry/.style={rectangle, rounded corners=2pt, draw=orange!80!black, fill=orange!14, minimum width=9mm, minimum height=6mm, inner sep=1pt}]

% (a) independent decoys
\begin{scope}
  \node[font=\small\bfseries, anchor=west] at (-0.2,2.25) {(a) independent law $\mathcal D_1$};
  \node[dec] (d1a) at (0,1.2) {$x_1$};
  \node[dec] (d2a) at (1.15,1.2) {$x_2$};
  \node at (2.15,1.2) {$\cdots$};
  \node[dec] (d3a) at (3.15,1.2) {$x_{N-1}$};
  \node[src] (sa) at (4.45,1.2) {$x_N$};
  \node[qry] (qa) at (2.2,-0.55) {$b\ \to\ y$};
  \foreach \d in {d1a,d2a,d3a} { \draw[->, black!45] (\d.south) -- (qa.north); }
  \draw[->, thick, royalazure] (sa.south) -- (qa.north);
  \node[font=\scriptsize, align=center, text=black!70] at (2.2,-1.45)
     {decoys independent of $x_N$:\quad $\E[y b\,x_j]=0$};
  \node[font=\scriptsize, align=center, text=royalazure] at (2.2,-2.05)
     {rule credit at uniform attention $=\frac{\gamma}{N}\cdot 1$};
\end{scope}

% (b) one anti-copy at a random position
\begin{scope}[shift={(7.2,0)}]
  \node[font=\small\bfseries, anchor=west] at (-0.2,2.25) {(b) anti-copy law $\mathcal D_0$};
  \node[dec] (d1b) at (0,1.2) {$x_1$};
  \node[anti] (d2b) at (1.15,1.2) {$-x_N$};
  \node at (2.15,1.2) {$\cdots$};
  \node[dec] (d3b) at (3.15,1.2) {$x_{N-1}$};
  \node[src] (sb) at (4.45,1.2) {$x_N$};
  \node[qry] (qb) at (2.2,-0.55) {$b\ \to\ y$};
  \foreach \d in {d1b,d3b} { \draw[->, black!45] (\d.south) -- (qb.north); }
  \draw[->, red!70!black] (d2b.south) -- (qb.north);
  \draw[->, thick, royalazure] (sb.south) -- (qb.north);
  \node[font=\scriptsize, text=red!70!black] at (1.15,1.75) {random position};
  \node[font=\scriptsize, align=center, text=black!70] at (2.2,-1.45)
     {each decoy: $\E[y b\,x_j]=-\frac{\gamma}{N-1}$;\quad multiset $\{x_j\}$ independent of $x_N$};
  \node[font=\scriptsize, align=center, text=red!70!black] at (2.2,-2.05)
     {rule credit at uniform attention $=\frac{\gamma}{N}\bigl(1-(N-1)\cdot\frac{1}{N-1}\bigr)=0$};
\end{scope}
\end{tikzpicture}
}
\caption{\textbf{Two laws with the same local statistics, read by uniform attention.} The query holds the current letter $b$ and
must predict the displayed successor $y=b x_N \xi$, where $\xi$ is an error independent of everything else. In (a) the prefix
positions are independent of the source. In (b) one prefix position, chosen uniformly at random, displays $-x_N$, so each position
correlates with the source by $-\tfrac{1}{N-1}$. This is too small to change any local statistic, but the $N-1$ positions are read
together and cancel the source's contribution (\cref{lem:credit}); the multiset of displayed values is even independent of $x_N$ (\cref{lem:matched}).}
\label{fig:cancel}
\end{figure}

\paragraph{Contributions.}
\begin{enumerate}[leftmargin=*,itemsep=2pt,topsep=2pt]
\item \textbf{Matched laws, separated training times} (\cref{lem:matched,thm:separation}). Laws with the same error rate, marginals and Bayes predictor,
at total-variation distance $O(N^{-1/2})$, give routing times $\Theta(\tfrac{N}{\gamma}\log(1+\varepsilon^{-2}))$ and $\Theta(\tfrac{N}{\gamma\varepsilon})$ under correlation loss,
with absolute constants. A conserved quantity solves the coupled attention--MLP flow exactly. Over-cancellation has an exact barrier: routing succeeds if and only if
$\varepsilon^2>-F(q_0)$. Cross-entropy keeps both orders for escaping the bottleneck.
\item \textbf{Distributional cancellation} (\cref{lem:matched}). Under the anti-copy law the label is independent of everything uniform attention reads,
so every non-attention parameter of any one-block network sees pure label noise at the start; only routing can break the plateau.
\item \textbf{A dose--response law for clean data} (\cref{cor:dilution}). A fraction $\lambda$ of traces without cancellation gives
$\Theta(\tfrac{N}{\gamma}\min\{\tfrac{1}{\varepsilon},\tfrac{1+\log(1+\lambda/\varepsilon)}{\lambda}\})$: the bottleneck disappears once $\lambda$ exceeds the initialization scale.
\item \textbf{Wrong CoT steps can cause it} (\cref{thm:causal}). A causal error kernel with local error rate $\eta$ and a constant fraction of wrong transitions
produces cancellation at every supervised step, with the correct step still Bayes optimal.
\end{enumerate}

\paragraph{The scope of the claim.} The theorems concern a one-block attention--ReLU network with specified matrix directions and four trained scalars,
derived exactly from its matrices (\cref{prop:realization}). The separation is in the initialization scale at fixed $N$; it becomes a separation in $N$ only if
$\varepsilon$ is tied to $N$. Complete cross-entropy routing is verified numerically rather than proved. In untied experiments every matrix entry trains:
the separation appears when training starts from the specified directions, with or without added noise, and does not appear from fully random initialization in our runs.
